\documentclass[a4paper,10pt]{article}
\usepackage[utf8]{inputenc}
\usepackage[T1]{fontenc}
\usepackage[french]{babel}
\usepackage{geometry}
\geometry{left=1.5cm,right=1.5cm,top=1.2cm,bottom=1.2cm}
\usepackage{amsmath,amssymb}
\usepackage{enumitem}
\usepackage{array}
\usepackage{booktabs}
\usepackage{xcolor}
\usepackage{tcolorbox}
\usepackage{tikz}
\usetikzlibrary{decorations.pathmorphing}
\usepackage{multicol}

\definecolor{vertpro}{HTML}{2da44e}
\definecolor{orange}{HTML}{d97706}
\definecolor{bleupro}{HTML}{0056b3}
\definecolor{rouge}{HTML}{c53030}
\definecolor{gris}{HTML}{718096}

\tcbuselibrary{skins,breakable}

\newtcolorbox{defbox}[1]{
  colback=green!3,
  colframe=vertpro,
  title={\small\textbf{#1}},
  fonttitle=\bfseries,
  boxsep=3pt,
  left=5pt,right=5pt,top=3pt,bottom=3pt
}

\newtcolorbox{methbox}[1]{
  colback=blue!3,
  colframe=bleupro,
  title={\small\textbf{#1}},
  fonttitle=\bfseries,
  boxsep=3pt,
  left=5pt,right=5pt,top=3pt,bottom=3pt
}

\newtcolorbox{attbox}{
  colback=red!3,
  colframe=rouge,
  title={\small\textbf{Attention !}},
  fonttitle=\bfseries,
  boxsep=3pt,
  left=5pt,right=5pt,top=3pt,bottom=3pt
}

\newtcolorbox{exbox}{
  colback=yellow!5,
  colframe=orange,
  boxsep=3pt,
  left=5pt,right=5pt,top=3pt,bottom=3pt
}

\setlength{\parindent}{0pt}
\setlength{\parskip}{4pt}
\pagestyle{empty}

\begin{document}

\begin{center}
{\Large\bfseries\color{vertpro} Fiche de révision — CCF Physique-Chimie}\\[4pt]
{\normalsize Terminale Bac Pro ERA-MA — Groupement 3 \quad | \quad Électricité \& Acoustique}
\end{center}

\vspace{6pt}
\begin{multicols}{2}

% ========================================
{\large\bfseries\color{vertpro} ÉLECTRICITÉ — TRANSPORT}
\vspace{4pt}

\begin{defbox}{Le réseau électrique}
\begin{center}
\begin{tikzpicture}[scale=0.5, every node/.style={font=\tiny}]
  \draw[fill=green!10, rounded corners] (0,0) rectangle (2,1);
  \node[align=center] at (1,0.5) {\textbf{Centrale}\\20 kV};
  \draw[->] (2,0.5) -- (3,0.5);
  \draw[fill=yellow!15, rounded corners] (3,0) rectangle (5.5,1);
  \node[align=center] at (4.25,0.5) {\textbf{Transfo}\\élévateur};
  \draw[->, rouge] (5.5,0.5) -- (7,0.5) node[midway,above]{\color{rouge}400 kV};
  \draw[fill=yellow!15, rounded corners] (7,0) rectangle (9.5,1);
  \node[align=center] at (8.25,0.5) {\textbf{Transfo}\\abaisseur};
  \draw[->] (9.5,0.5) -- (11,0.5);
  \draw[fill=blue!5, rounded corners] (11,0) rectangle (13.5,1);
  \node[align=center] at (12.25,0.5) {\textbf{Atelier}\\230/400 V};
\end{tikzpicture}
\end{center}
\end{defbox}

\begin{defbox}{Formules essentielles}
\renewcommand{\arraystretch}{1.4}
\begin{tabular}{>{\bfseries\small}l l}
Puissance & $P = U \times I$ \\
Loi d'Ohm & $U = R \times I$ \\
Pertes Joule & $\boxed{P_J = R \times I^2}$ \\
Énergie & $E = P \times t$ \\
Résistance câble & $R = \rho \times \dfrac{L}{S}$ \\[6pt]
Transformateur & $\dfrac{U_2}{U_1} = \dfrac{N_2}{N_1}$ \\
\end{tabular}
\end{defbox}

\begin{defbox}{Le transformateur}
\begin{itemize}[nosep,leftmargin=12pt]
  \item \textbf{Élévateur} : $U_2 > U_1$ (monte la tension)
  \item \textbf{Abaisseur} : $U_2 < U_1$ (baisse la tension)
  \item Le rapport $N_2/N_1$ détermine le type
  \item Rôle : adapter la tension sans changer la fréquence
\end{itemize}
\end{defbox}

\begin{defbox}{Pertes par effet Joule}
Quand le courant traverse un câble de résistance $R$ :
$$P_J = R \times I^2$$
Les pertes chauffent les câbles $\to$ énergie perdue.

\textbf{Comment les réduire ?}
\begin{itemize}[nosep,leftmargin=12pt]
  \item Réduire $R$ (câbles plus gros, cuivre)
  \item \textbf{Réduire $I$} en augmentant $U$ (haute tension)
\end{itemize}
\end{defbox}

\begin{methbox}{Pourquoi transporter en haute tension ?}
À puissance constante : $I = \dfrac{P}{U}$

Si $U$ est multiplié par 10 $\to$ $I$ est divisé par 10 $\to$ $P_J = RI^2$ est divisé par \textbf{100} !

C'est pour ça que le réseau transporte en 400\,000 V.
\end{methbox}

\begin{exbox}
\textbf{Exemple :} $P = 15$ kW, $R = 0{,}1\;\Omega$\\[4pt]
En 400 V : $I = 37{,}5$ A $\to$ $P_J = 0{,}1 \times 37{,}5^2 = 140{,}6$ W\\
En 20\,000 V : $I = 0{,}75$ A $\to$ $P_J = 0{,}1 \times 0{,}75^2 = 0{,}056$ W\\[4pt]
$\to$ \textbf{2\,500 fois} moins de pertes en haute tension !
\end{exbox}

\begin{attbox}
\begin{itemize}[nosep,leftmargin=12pt]
  \item Convertir kW en W avant de calculer ($15$ kW $= 15\,000$ W)
  \item $I^2$ : ne pas oublier de mettre au carré !
  \item Unités : $P$ en W, $U$ en V, $I$ en A, $R$ en $\Omega$
\end{itemize}
\end{attbox}

\columnbreak

% ========================================
{\large\bfseries\color{vertpro} ACOUSTIQUE — ATTÉNUATION}
\vspace{4pt}

\begin{defbox}{Niveau sonore en décibels}
$$L = 10 \log\!\left(\frac{I}{I_0}\right) \quad \text{en dB}$$
$I_0 = 10^{-12}$ W/m² (seuil d'audibilité)

\begin{center}
\begin{tikzpicture}[scale=0.42]
  \draw[fill=green!15] (0,0) rectangle (8,1.2);
  \draw[fill=yellow!20] (8,0) rectangle (11.5,1.2);
  \draw[fill=red!15] (11.5,0) rectangle (15,1.2);
  \draw[thick, orange] (8,0) -- (8,1.2);
  \draw[thick, rouge] (11.5,0) -- (11.5,1.2);
  \node[font=\tiny] at (4,0.6) {Confort};
  \node[font=\tiny] at (9.75,0.6) {Danger};
  \node[font=\tiny] at (13.25,0.6) {Douleur};
  \node[below, font=\tiny] at (0,-0.1) {0};
  \node[below, font=\tiny] at (4,-0.1) {40};
  \node[below, font=\tiny\bfseries, orange] at (8,-0.1) {85};
  \node[below, font=\tiny\bfseries, rouge] at (11.5,-0.1) {120};
  \node[below, font=\tiny] at (15,-0.1) {140 dB};
\end{tikzpicture}
\end{center}
\end{defbox}

\begin{defbox}{Seuils réglementaires}
\renewcommand{\arraystretch}{1.3}
\begin{tabular}{|l|l|}
\hline
\textbf{85 dB} & Port de protections \textbf{obligatoire} \\
\hline
\textbf{120 dB} & Seuil de \textbf{douleur} \\
\hline
\textbf{50 dB} & Max recommandé (bureau) \\
\hline
\end{tabular}
\end{defbox}

\begin{defbox}{Affaiblissement acoustique $R$}
Quand le son traverse une paroi :
$$\boxed{L_{\text{transmis}} = L_{\text{incident}} - R}$$
$R$ en dB = indice d'affaiblissement de la paroi.

Plus $R$ est grand, meilleure est l'isolation.
\end{defbox}

\begin{defbox}{Isolation $\neq$ Absorption}
\renewcommand{\arraystretch}{1.4}
\begin{tabular}{|>{\small}p{3.2cm}|>{\small}p{3.2cm}|}
\hline
\textbf{ISOLATION} & \textbf{ABSORPTION} \\
\hline
Empêche le son de \textbf{traverser} une paroi & Réduit l'\textbf{écho} à l'intérieur d'une pièce \\
\hline
Masse, double paroi, désolidarisation & Matériaux poreux (laine, mousse, tissu) \\
\hline
Entre deux pièces & Dans une même pièce \\
\hline
\end{tabular}
\end{defbox}

\begin{methbox}{Méthode — Exercice type acoustique}
\begin{enumerate}[nosep,leftmargin=12pt]
  \item Lire $L_{\text{incident}}$ et $R$ de la cloison
  \item Calculer $L_{\text{transmis}} = L_{\text{incident}} - R$
  \item Comparer à la réglementation (50 dB bureau)
  \item Si insuffisant : augmenter $R$ (ajouter isolant)
  \item Isolation $\neq$ Absorption (ne pas confondre !)
\end{enumerate}
\end{methbox}

\begin{exbox}
\textbf{Exemple :} Atelier 95 dB, cloison $R = 40$ dB\\
$L_{\text{bureau}} = 95 - 40 = 55$ dB $> 50$ dB $\to$ \textbf{non conforme}\\[4pt]
Ajout isolant $+10$ dB : $L' = 95 - 50 = 45$ dB $\to$ \textbf{OK}
\end{exbox}

\begin{attbox}
\begin{itemize}[nosep,leftmargin=12pt]
  \item Laine de roche sur les murs $=$ absorption (pas isolation)
  \item L'isolation agit sur la cloison séparative
  \item Doubler l'épaisseur $\to$ $R$ augmente de seulement 5-6 dB
\end{itemize}
\end{attbox}

\end{multicols}

\end{document}
